\writeSectionFrame{\presentationTitle}{Übung}

\begin{frame}
  \titlepage
\end{frame}

\begin{frame}
  \frametitle{Ziel}
  \begin{itemize}
    \item Verstehen und Implementieren eines Connection-Pools, der das Multiton-Muster verwendet.
    \item Anwendung des Multiton-Musters zur Verwaltung einer begrenzten Anzahl von Instanzen.
  \end{itemize}
\end{frame}

\begin{frame}
  \frametitle{Aufgabenbeschreibung}
    \textbf{Pool-Klasse erstellen}
    \begin{itemize}
      \item Implementieren Sie eine Klasse \texttt{ConnectionPool}, die Verbindungen verwaltet.
      \item Die \texttt{ConnectionPool}-Klasse soll folgende Methoden bieten:
      \begin{itemize}
        \item \texttt{getConnection(String dbName)}: Gibt eine Verbindung für den angegebenen Datenbanknamen zurück, wenn eine verfügbar ist, oder erstellt eine neue, wenn der Pool noch nicht voll ist.
        \item \texttt{releaseConnection(String dbName)}: Gibt eine Verbindung für den angegebenen Datenbanknamen zurück, um sie im Pool wieder verfügbar zu machen.
        \item \texttt{printConnections()}: Zeigt alle Verbindungen im Pool an.
      \end{itemize}
    \end{itemize}
\end{frame}

\begin{frame}
  \frametitle{Aufgabenbeschreibung}
    \textbf{DatabaseConnection-Klasse erstellen}
    \begin{itemize}
     \item \texttt{DatabaseConnection}-Klasse braucht nur ein Attribut für den Datenbanknamen und eine toString()-Methode.
      \item Die \texttt{DatabaseConnection}-Instanz soll bei der Erstellung in den Pool eingefügt werden und beim Freigeben wieder in den Pool zurückgegeben werden.
      \end{itemize}
\end{frame}

\begin{frame}
  \frametitle{Aufgabenbeschreibung}
    \textbf{Verwenden des Pools}
    \begin{itemize}
      \item Implementieren Sie eine \texttt{main}-Methode oder eine Testklasse, die mehrere Verbindungen anfordert, zurückgibt und die Verwaltung des Pools testet.
    \end{itemize}
\end{frame}